% Chunk: tail of \S6.2, \S6.3 Momentum Space Rules, and opening of \S6.4.
% Source: Weinberg QFT Vol. I, physical pp. 303--309 (printed pp. 280--286).
% Status: source-reviewed, compile-clean, and visually checked; equations (6.2.30)--(6.3.13); Figure 6.9 integrated with caption; one footnote.

Also, $\psi^{(-)\dagger}$ and $\psi^{(+)}$ would annihilate a vacuum state on
the right, and $\psi^{(-)}$ and $\psi^{(+)\dagger}$ would annihilate the vacuum
state on the left, so $\psi^{(+)}$ and $\psi^{(-)}$ may be replaced everywhere
in Eq. \eqref{eq:6.2.29} with the complete field $\psi=\psi^{(+)}+\psi^{(-)}$:
\begin{equation}
-i\Delta_{\ell m}(x,y)=\theta(x-y)(\psi_\ell(x)\psi_m^\dagger(y))_0
\mathbin{\pm}\theta(y-x)(\psi_m^\dagger(y)\psi_\ell(x))_0.
\tag{6.2.30}
\label{eq:6.2.30}
\end{equation}
This is often written
\begin{equation}
-i\Delta_{\ell m}(x,y)=(T\{\psi_\ell(x)\psi_m^\dagger(y)\})_0,
\tag{6.2.31}
\label{eq:6.2.31}
\end{equation}
where $T$ is a time-ordered product, whose definition is now extended
to all fields, with a minus sign for any odd permutation of fermionic
operators.\footnote{This is not inconsistent with our previous definition of
the time-ordered product of Hamiltonian densities in Chapter 3, because the
Hamiltonian density can only contain even numbers of fermionic field factors.}

\subsection{Momentum Space Rules}

The Feynman rules outlined in Section 6.1 specify how to calculate the
contribution to the $S$-matrix of a given $N$th order diagram, as the integral
over $N$ spacetime coordinates of a product of spacetime-dependent factors.
For a final particle (or antiparticle) line with momentum $p'^\mu$ leaving a
vertex with spacetime coordinate $x^\mu$, we get a factor proportional to
$\exp(-ip'\cdot x)$, and for an initial particle line with momentum $p^\mu$
entering a vertex with spacetime coordinate $x^\mu$, we get a factor
proportional to $\exp(+ip\cdot x)$. In Section 6.2 we saw that the factor
associated with an internal line running from $y$ to $x$ can be expressed as a
Fourier integral, over off-shell four-momenta $q^\mu$, of an integrand
proportional to $\exp(iq\cdot(x-y))$. We can think of $q^\mu$ as the
four-momentum flowing along the internal line in the direction of the arrow
from $y$ to $x$. Hence the integral over each vertex's spacetime position
merely yields a factor
\begin{equation}
(2\pi)^4\delta^4\left(\sum p+\sum q-\sum p'-\sum q'\right),
\tag{6.3.1}
\label{eq:6.3.1}
\end{equation}
where $\sum p'$ and $\sum p$ denote the total four-momentum of all the final or
initial particles leaving or entering the vertex, and $\sum q'$ and $\sum q$
denote the total four-momentum of all the internal lines with arrows leaving or
entering the vertex, respectively. Of course, in place of these integrals over
$x^\mu$s, we now have to do integrals over the Fourier variables $q^\mu$, one
for each internal line.

These considerations can be encapsulated in a new set of Feynman rules

\begingroup
\renewcommand{\thefigure}{6.\arabic{figure}}
\setcounter{figure}{8}
\begin{figure}[htbp]
% Figure 6.9 plate only; caption remains in sec63.tex.
\begin{center}
\begin{tikzpicture}[
  x=1cm,
  y=1cm,
  line cap=round,
  line join=round,
  pairing/.style={draw, line width=0.7pt},
  every node/.style={font=\small}
]
  % (a) Final particle line.
  \draw[pairing] (-5.7,3.75) -- (-5.7,6.25);
  \draw[-{Latex[length=2.3mm,width=1.5mm]},line width=0.7pt]
    (-5.7,4.75) -- (-5.7,5.28);
  \fill (-5.7,3.75) circle (1.35pt);
  \draw[pairing] (-5.7,3.75) -- (-5.90,3.43);
  \draw[pairing] (-5.7,3.75) -- (-5.50,3.43);
  \node[above] at (-5.7,6.25) {$\mathbf p'\sigma'n'$};
  \node[right] at (-5.50,3.78) {$\ell$};
  \node[anchor=west] at (-5.0,5.05)
    {$(2\pi)^{-3/2}u_\ell^*(\mathbf p'\sigma'n')$};
  \node at (-4.35,3.13) {(a)};

  % (b) Final antiparticle line.
  \draw[pairing] (-1.2,3.75) -- (-1.2,6.25);
  \draw[-{Latex[length=2.3mm,width=1.5mm]},line width=0.7pt]
    (-1.2,5.28) -- (-1.2,4.75);
  \fill (-1.2,3.75) circle (1.35pt);
  \draw[pairing] (-1.2,3.75) -- (-1.40,3.43);
  \draw[pairing] (-1.2,3.75) -- (-1.00,3.43);
  \node[above] at (-1.2,6.25) {$\mathbf p'\sigma'n'^{c}$};
  \node[right] at (-1.00,3.78) {$\ell$};
  \node[anchor=west] at (-0.5,5.05)
    {$(2\pi)^{-3/2}v_\ell(\mathbf p'\sigma'n')$};
  \node at (0.15,3.13) {(b)};

  % (c) Initial particle line.
  \draw[pairing] (3.3,3.75) -- (3.3,6.25);
  \draw[-{Latex[length=2.3mm,width=1.5mm]},line width=0.7pt]
    (3.3,4.75) -- (3.3,5.28);
  \fill (3.3,6.25) circle (1.35pt);
  \draw[pairing] (3.3,6.25) -- (3.10,6.57);
  \draw[pairing] (3.3,6.25) -- (3.50,6.57);
  \node[below] at (3.3,3.75) {$\mathbf p\sigma n$};
  \node[right] at (3.50,6.23) {$\ell$};
  \node[anchor=west] at (4.0,5.05)
    {$(2\pi)^{-3/2}u_\ell(\mathbf p\sigma n)$};
  \node at (4.65,3.13) {(c)};

  % (d) Initial antiparticle line.
  \draw[pairing] (-5.7,-0.45) -- (-5.7,2.05);
  \draw[-{Latex[length=2.3mm,width=1.5mm]},line width=0.7pt]
    (-5.7,1.08) -- (-5.7,0.55);
  \fill (-5.7,2.05) circle (1.35pt);
  \draw[pairing] (-5.7,2.05) -- (-5.90,2.37);
  \draw[pairing] (-5.7,2.05) -- (-5.50,2.37);
  \node[below] at (-5.7,-0.45) {$\mathbf p\sigma n^{c}$};
  \node[right] at (-5.50,2.03) {$\ell$};
  \node[anchor=west] at (-5.0,0.83)
    {$(2\pi)^{-3/2}v_\ell^*(\mathbf p\sigma n)$};
  \node at (-4.35,-0.92) {(d)};

  % (e) Direct initial-final pairing, particle or antiparticle.
  \draw[pairing] (-0.1,-0.45) -- (-0.1,2.05);
  \draw[-{Latex[length=2.3mm,width=1.5mm]},line width=0.7pt]
    (-0.1,0.55) -- (-0.1,1.08);
  \node[above] at (-0.1,2.05) {$\mathbf p'\sigma'n'$};
  \node[below] at (-0.1,-0.45) {$\mathbf p\sigma n$};

  \node at (0.80,0.80) {or};

  \draw[pairing] (1.7,-0.45) -- (1.7,2.05);
  \draw[-{Latex[length=2.3mm,width=1.5mm]},line width=0.7pt]
    (1.7,1.08) -- (1.7,0.55);
  \node[above] at (1.7,2.05) {$\mathbf p'\sigma'n'^{c}$};
  \node[below] at (1.7,-0.45) {$\mathbf p\sigma n^{c}$};
  \node[anchor=west] at (2.55,0.83)
    {$\delta^3(\mathbf p'-\mathbf p)\delta_{\sigma'\sigma}\delta_{n'n}$};
  \node at (0.80,-0.92) {(e)};

  % (f) Internal line, with momentum flowing from m to ell.
  \draw[pairing] (-2.20,-2.35) -- (0.75,-2.35);
  \draw[-{Latex[length=2.3mm,width=1.5mm]},line width=0.7pt]
    (-0.55,-2.35) -- (-1.08,-2.35);
  \fill (-2.20,-2.35) circle (1.25pt);
  \fill (0.75,-2.35) circle (1.25pt);
  \draw[pairing] (-2.20,-2.35) -- (-2.52,-2.17);
  \draw[pairing] (-2.20,-2.35) -- (-2.52,-2.53);
  \draw[pairing] (0.75,-2.35) -- (1.07,-2.17);
  \draw[pairing] (0.75,-2.35) -- (1.07,-2.53);
  \node[above] at (-0.72,-2.35) {$q$};
  \node[below] at (-2.05,-2.35) {$\ell$};
  \node[below] at (0.60,-2.35) {$m$};
  \node[anchor=west] at (1.55,-2.35)
    {$\displaystyle \frac{-i}{(2\pi)^4}
      \frac{P_{\ell m}(q)}{q^2+m_\ell^2-i\epsilon}$};
  \node at (-0.72,-3.15) {(f)};
\end{tikzpicture}
\end{center}

\caption{Graphical representation of pairings of operators arising in the
momentum-space evaluation of the $S$-matrix. The expressions on the right are
the factors that must be included in the momentum space integrand of the
$S$-matrix for each line of the Feynman diagram.}
\label{fig:6.9}
\end{figure}
\endgroup

(see Figure~\ref{fig:6.9}) for calculating contribution to the $S$-matrix as integrals
over momentum variables:

\begin{enumerate}
\item[(i)] Draw all Feynman diagrams of the desired order, just as described
in Section 6.1. However, instead of labelling each vertex with a spacetime
coordinate, each internal line is now labelled with an off-mass-shell
four-momentum, considered conventionally to flow in the direction of the arrow
(or in either direction for neutral particle lines without arrows.)

\item[(ii)] For each vertex of type $j$, include a factor
\begin{equation}
-i(2\pi)^4g_j\delta^4\left(\sum p+\sum q-\sum p'-\sum q'\right)
\tag{6.3.2}
\label{eq:6.3.2}
\end{equation}
with the momentum sums having the same meaning as in \eqref{eq:6.3.1}. This delta
function ensures that the four-momentum is conserved at every point in the
diagram. For each external line running upwards out of the diagram, include a
factor
\[
(2\pi)^{-3/2}u_\ell^*(\mathbf p'\sigma'n')
\quad\hbox{or}\quad
(2\pi)^{-3/2}v_\ell(\mathbf p'\sigma'n'),
\]
for arrows pointing upwards or downwards, respectively. For each external line
running from below into the diagram, include a factor
\[
(2\pi)^{-3/2}u_\ell(\mathbf p\sigma n)
\quad\hbox{or}\quad
(2\pi)^{-3/2}v_\ell^*(\mathbf p\sigma n),
\]
for arrows pointing upwards or downwards, respectively. For each internal line
with ends labelled $\ell$ and
$m$, the arrow pointing from $m$ to $\ell$, and carrying a momentum label
$q^\mu$, include as a factor the integrand of the integral for
$-i\Delta_{\ell m}(q)$:
\begin{equation}
-i(2\pi)^{-4}\frac{P_{\ell m}(q)}{q^2+m_\ell^2-i\epsilon}.
\tag{6.3.3}
\label{eq:6.3.3}
\end{equation}
A reminder: for scalars or antiscalars of four-momentum $q$, the $u$s and $v$s
are simply $(2q^0)^{-1/2}$, while the polynomial $P(q)$ is unity. For Dirac
spinors of four-momentum $p$ and mass $M$, the $u$s and $v$s are the normalized
Dirac spinors described in Section 5.5, and the polynomial $P(p)$ is the matrix
$(-i\gamma_\mu p^\mu+M)\beta$.

\item[(iii)] Integrate the product of all these factors over the four-momenta
carried by internal lines, and sum over all field indices $\ell,m$, etc.

\item[(iv)] Add up the results obtained in this way from each Feynman diagram.
\end{enumerate}

Additional combinatoric factors and fermionic signs may need to be included,
as described in parts (v) and (vi) of Section 6.1. Examples will be given at
the end of this section.

We have a four-momentum integration variable for every internal line, but many
of these are eliminated by the delta functions associated with vertices. Since
energy and momentum are separately conserved for each connected part of a
Feynman diagram, there will be $C$ delta functions left over in a graph with
$C$ connected parts. Hence in a diagram with $I$ internal lines and $V$
vertices, the number of independent four-momenta that are not fixed by the
delta functions is $I-[V-C]$. This is clearly also the number $L$ of
independent loops:
\begin{equation}
L=I-V+C,
\tag{6.3.4}
\label{eq:6.3.4}
\end{equation}
which is defined as the maximum number of internal lines that can be cut
without disconnecting the diagram, because any such and only such internal
lines can be assigned an independent four-momentum. We can think of the
independent momentum variables as characterizing the momenta that circulate in
each loop. In particular a \textit{tree} graph is one without loops; after
taking the delta functions into account there are no momentum-space integrals
left for such graphs.

For instance, in a theory with interaction \eqref{eq:6.1.18}, the $S$-matrix \eqref{eq:6.1.27}
for fermion--boson scattering is given by the momentum-space Feynman rules as
\begin{align*}
&S_{\mathbf p'_1\sigma'_1n'_1\,\mathbf p'_2\sigma'_2n'_2,\,
\mathbf p_1\sigma_1n_1\,\mathbf p_2\sigma_2n_2}\\
&=\sum_{k'\ell'm'k\ell m}(-i)^2(2\pi)^8g_{\ell'm'k'}g_{m\ell k}
u_{\ell'}^*(\mathbf p'_1\sigma'_1n'_1)u_\ell(\mathbf p_1\sigma_1n_1)\\
&\quad\cdot\int d^4q\left(-i(2\pi)^{-4}
\frac{P_{m'm}(q)}{q^2+m_m^2-i\epsilon}\right)(2\pi)^{-6}\\
&\quad\cdot\Bigl[
u_{k'}^*(\mathbf p'_2\sigma'_2n'_2)u_k(\mathbf p_2\sigma_2n_2)
\delta^4(p_2-p'_2-q)\delta^4(p_1+q-p'_1)\\
&\qquad{}+u_k^*(\mathbf p'_2\sigma'_2n'_2)u_{k'}(\mathbf p_2\sigma_2n_2)
\delta^4(p_2-p'_1-q)\delta^4(p_1-p'_2+q)\Bigr],
\end{align*}
with labels 1 and 2 here denoting fermions and bosons, respectively. The
momentum-space integral here is trivial, and gives
\begin{align}
&S_{\mathbf p'_1\sigma'_1n'_1\,\mathbf p'_2\sigma'_2n'_2,\,
\mathbf p_1\sigma_1n_1\,\mathbf p_2\sigma_2n_2}\notag\\
&=i(2\pi)^{-2}\delta^4(p_1+p_2-p'_1-p'_2)
\sum_{k'\ell'm'k\ell m}g_{\ell'm'k'}g_{m\ell k}
u_{\ell'}^*(\mathbf p'_1\sigma'_1n'_1)u_\ell(\mathbf p_1\sigma_1n_1)
\notag\\
&\quad\cdot\left[
\frac{P_{m'm}(p'_1-p_1)}{(p'_1-p_1)^2+m_m^2-i\epsilon}
u_{k'}^*(\mathbf p'_2\sigma'_2n'_2)u_k(\mathbf p_2\sigma_2n_2)\right.
\notag\\
&\qquad\left.{}+
\frac{P_{m'm}(p'_2-p_1)}{(p'_2-p_1)^2+m_m^2-i\epsilon}
u_k^*(\mathbf p'_2\sigma'_2n'_2)u_{k'}(\mathbf p_2\sigma_2n_2)
\right].
\tag{6.3.5}
\label{eq:6.3.5}
\end{align}
In the same way, the $S$-matrix element \eqref{eq:6.1.28} for fermion--fermion
scattering in the same theory is
\begin{align}
&S_{\mathbf p'_1\sigma'_1n'_1\,\mathbf p'_2\sigma'_2n'_2,\,
\mathbf p_1\sigma_1n_1\,\mathbf p_2\sigma_2n_2}\notag\\
&=i(2\pi)^{-2}\delta^4(p_1+p_2-p'_1-p'_2)
\sum_{k'\ell'm'k\ell m}g_{\ell m'k'}g_{\ell'mk}
\frac{P_{k'k}(p'_1-p_1)}{(p'_1-p_1)^2+m_k^2-i\epsilon}
\notag\\
&\quad\cdot u_\ell^*(\mathbf p'_2\sigma'_2n'_2)
u_{\ell'}^*(\mathbf p'_1\sigma'_1n'_1)
u_{m'}(\mathbf p_2\sigma_2n_2)u_m(\mathbf p_1\sigma_1n_1)
-[1'\rightleftarrows2'].
\tag{6.3.6}
\label{eq:6.3.6}
\end{align}
These results illustrate the need for a more compact notation. We may define a
fermion--boson coupling matrix
\begin{equation}
[\Gamma_k]_{\ell m}\equiv g_{\ell mk}.
\tag{6.3.7}
\label{eq:6.3.7}
\end{equation}
The matrix elements \eqref{eq:6.3.5} and \eqref{eq:6.3.6} for fermion--boson and
fermion--fermion scattering can then be rewritten in matrix notation as
\begin{align}
&S_{\mathbf p'_1\sigma'_1n'_1\,\mathbf p'_2\sigma'_2n'_2,\,
\mathbf p_1\sigma_1n_1\,\mathbf p_2\sigma_2n_2}\notag\\
&=i(2\pi)^{-2}\delta^4(p_1+p_2-p'_1-p'_2)\sum_{k'k}
\Biggl[\left(u^\dagger(\mathbf p'_1\sigma'_1n'_1)\Gamma_{k'}
\frac{P(p'_1-p_1)}{(p'_1-p_1)^2+M^2-i\epsilon}
\Gamma_k u(\mathbf p_1\sigma_1n_1)\right)\notag\\
&\qquad\cdot u_{k'}^*(\mathbf p'_2\sigma'_2n'_2)
u_k(\mathbf p_2\sigma_2n_2)\notag\\
&\quad{}+\left(u^\dagger(\mathbf p'_1\sigma'_1n'_1)\Gamma_{k'}
\frac{P(p'_2-p_1)}{(p'_2-p_1)^2+M^2-i\epsilon}
\Gamma_k u(\mathbf p_1\sigma_1n_1)\right)\notag\\
&\qquad\cdot u_k^*(\mathbf p'_2\sigma'_2n'_2)
u_{k'}(\mathbf p_2\sigma_2n_2)\Biggr].
\tag{6.3.8}
\label{eq:6.3.8}
\end{align}
and
\begin{align}
&S_{\mathbf p'_1\sigma'_1n'_1\,\mathbf p'_2\sigma'_2n'_2,\,
\mathbf p_1\sigma_1n_1\,\mathbf p_2\sigma_2n_2}\notag\\
&=i(2\pi)^{-2}\delta^4(p_1+p_2-p'_1-p'_2)
\sum_{k'k}\frac{P_{k'k}(p'_1-p_1)}{(p'_1-p_1)^2+m_k^2-i\epsilon}\notag\\
&\quad\cdot
\left(u^\dagger(\mathbf p'_2\sigma'_2n'_2)\Gamma_{k'}
u(\mathbf p'_1\sigma'_1n'_1)\right)
\left(u^\dagger(\mathbf p_2\sigma_2n_2)\Gamma_k
u(\mathbf p_1\sigma_1n_1)\right)-[1'\rightleftarrows2'].
\tag{6.3.9}
\label{eq:6.3.9}
\end{align}
where $M^2$ and $m^2$ are the diagonal mass matrices of the fermions and bosons
in Eqs. \eqref{eq:6.3.8} and \eqref{eq:6.3.9}, respectively. The general rule is that in using
matrix notation, one writes coefficient functions, coupling matrices, and
propagators in an order dictated by following lines \textit{backwards} from the
order indicated by the arrows. In the same notation, the $S$-matrix for
boson--boson scattering in the same theory would be given by a sum of one-loop
diagrams, shown in Figure~\ref{fig:6.7}:
\begin{align}
&S_{\mathbf p'_1\sigma'_1n'_1\,\mathbf p'_2\sigma'_2n'_2,\,
\mathbf p_1\sigma_1n_1\,\mathbf p_2\sigma_2n_2}\notag\\
&=-\delta^4(p_1+p_2-p'_1-p'_2)
\sum_{k_1k_2k'_1k'_2}
u_{k'_1}^*(\mathbf p'_1\sigma'_1n'_1)
u_{k'_2}^*(\mathbf p'_2\sigma'_2n'_2)
u_{k_1}(\mathbf p_1\sigma_1n_1)u_{k_2}(\mathbf p_2\sigma_2n_2)\notag\\
&\quad\cdot\int d^4q\,\operatorname{Tr}\left\{
\Gamma_{k'_2}\frac{P(q)}{q^2+M^2-i\epsilon}
\Gamma_{k'_1}\frac{P(q+p'_1)}{(q+p'_1)^2+M^2-i\epsilon}\right.\notag\\
&\qquad\left.\cdot
\Gamma_{k_1}\frac{P(q+p'_1-p_1)}{(q+p'_1-p_1)^2+M^2-i\epsilon}
\Gamma_{k_2}\frac{P(q-p'_2)}{(q-p'_2)^2+M^2-i\epsilon}
\right\}+\cdots.
\tag{6.3.10}
\label{eq:6.3.10}
\end{align}
where the ellipsis in the last line indicates terms obtained by permuting
bosons $1'$, $2'$, and $2$. The minus sign at the beginning of the right-hand
side is the extra minus sign associated with fermionic loops. Note that after
elimination of delta functions there is just one momentum-space integral here,
as appropriate for a diagram with one loop. We shall see how to do this sort of
momentum-space integral in Chapter 11.

To make this more specific, consider a theory with a Dirac spinor field
$\psi(x)$ of mass $M$ and a pseudoscalar field $\phi(x)$ of mass $m$,
interacting through the interaction $-ig\bar\psi\gamma_5\psi\phi$. (The factor
$-i$ is inserted to make this interaction Hermitian for real coupling
constants $g$.) Recall that the polynomial $P(q)$ for the scalar is just unity,
while for the spinor it is $[-i\gamma_\mu q^\mu+M]\beta$. Also, the $u$ for a
scalar of energy $E$ is $(2E)^{-1/2}$, while for the spinor $u$ is the
conventionally normalized Dirac spinor discussed in Section 5.5. Eqs. \eqref{eq:6.3.8},
\eqref{eq:6.3.9}, and \eqref{eq:6.3.10} give the lowest-order connected $S$-matrix elements for
fermion--boson scattering, fermion--fermion scattering, and boson--boson
scattering:
\begin{align*}
S_{\mathbf p'_1\sigma'_1\,\mathbf p'_2,\,\mathbf p_1\sigma_1\,\mathbf p_2}
={}&-i(2\pi)^{-2}g^2(4E'_2E_2)^{-1/2}
\delta^4(p_1+p_2-p'_1-p'_2)\\
&\cdot\Biggl[
\bar u(\mathbf p'_1\sigma'_1)\gamma_5
\frac{-i\gamma_\mu(p'_1-p_1)^\mu+M}{(p'_1-p_1)^2+m^2-i\epsilon}
\gamma_5u(\mathbf p_1\sigma_1)\\
&\qquad{}+\bar u(\mathbf p'_1\sigma'_1)\gamma_5
\frac{-i\gamma_\mu(p'_2-p_1)^\mu+M}{(p'_2-p_1)^2+m^2-i\epsilon}
\gamma_5u(\mathbf p_1\sigma_1)\Biggr],
\\[4pt]
S_{\mathbf p'_1\sigma'_1\,\mathbf p'_2\sigma'_2,\,
\mathbf p_1\sigma_1\,\mathbf p_2\sigma_2}
={}&i(2\pi)^{-2}g^2\delta^4(p_1+p_2-p'_1-p'_2)\\
&\cdot\bigl(\bar u(\mathbf p'_2\sigma'_2)\gamma_5u(\mathbf p'_1\sigma'_1)\bigr)
\bigl(\bar u(\mathbf p_2\sigma_2)\gamma_5u(\mathbf p_1\sigma_1)\bigr)\\
&\cdot\frac{1}{(p'_1-p_1)^2+m^2-i\epsilon}-[1'\rightleftarrows2'],
\\[4pt]
S_{\mathbf p'_1\,\mathbf p'_2,\,\mathbf p_1\,\mathbf p_2}
={}&-(2\pi)^{-6}g^2(16E_1E_2E'_1E'_2)^{-1/2}
\delta^4(p_1+p_2-p'_1-p'_2)\\
&\cdot\int d^4q\,\operatorname{Tr}\Biggl\{
\gamma_5\frac{-i\gamma_\mu q^\mu+M}{q^2+M^2-i\epsilon}
\gamma_5\frac{-i\gamma_\mu(q+p'_1)^\mu+M}{(q+p'_1)^2+M^2-i\epsilon}\\
&\qquad\cdot\gamma_5
\frac{-i\gamma_\mu(q+p'_1-p_1)^\mu+M}{(q+p'_1-p_1)^2+M^2-i\epsilon}
\gamma_5\frac{-i\gamma_\mu(q-p'_2)^\mu+M}{(q-p'_2)^2+M^2-i\epsilon}
\Biggr\}+\cdots,
\end{align*}
where in the last formula the ellipsis indicates a sum over permutations of
particles $2$, $1'$, $2'$. The factors $\beta$ in the fermion propagator
numerators have been used to replace $u^\dagger$ with $\bar u$.

\begin{center}
$*\quad *\quad *$
\end{center}

Another useful topological result expresses a sort of conservation law of
lines. For the moment we can think of all internal and external lines as being
created at vertices and destroyed in pairs at the centers of internal lines or
when external lines leave the diagram. (This has nothing to do with the
directions of the arrows carried by these lines.) Equating the numbers of lines
that are created and destroyed then gives
\begin{equation}
2I+E=\sum_i n_iV_i,
\tag{6.3.11}
\label{eq:6.3.11}
\end{equation}
where $I$ and $E$ are the numbers of internal and external lines, $V_i$ are the
numbers of vertices of various types labelled $i$, and $n_i$ are the number of
lines attached to each vertex. (This also holds separately for fields of each
type.) In particular, if all interactions involve the same number $n_i=n$ of
fields, then this reads
\begin{equation}
2I+E=nV,
\tag{6.3.12}
\label{eq:6.3.12}
\end{equation}
where $V$ is the total number of all vertices. In this case, we can eliminate
$I$ from Eqs. \eqref{eq:6.3.4} and \eqref{eq:6.3.11}, and find that for a connected (i.e.,
$C=1$) graph the number of vertices is given by
\begin{equation}
V=\frac{2L+E-2}{n-2}.
\tag{6.3.13}
\label{eq:6.3.13}
\end{equation}
For instance for a trilinear interaction the diagrams for a scattering process
($E=4$) with $L=0,1,2\cdots$ has $V=2,4,6\cdots$ vertices. In general, the
expansion in powers of the coupling constants is an expansion in increasing
numbers of loops.

\subsection{Off the Mass Shell}

In the Feynman diagrams for any $S$-matrix element all external lines are ``on
the mass shell''; that is, the four-momentum associated with an external line
for a particle of mass $m$ is constrained to satisfy $p_\mu p^\mu=-m^2$. It is
often important also to consider Feynman diagrams ``off the mass shell'', for
which the external line energies like the energies associated with internal
lines are free variables, unrelated to any three-momenta. For one thing, these
arise as parts of larger Feynman diagrams; for instance, a loop appearing as an
insertion in some internal line of a diagram could be regarded as a Feynman
diagram with two external lines, both off the mass shell.

Of course, once we calculate the contribution of a given Feynman diagram off
the mass shell, it is easy to calculate the associated $S$-matrix elements by
going to the mass shell, taking the four-momentum $p^\mu$ flowing along the line
\textit{into} the diagram to have $p^0=\sqrt{\mathbf p^2+m^2}$ for particles in
the initial state and $p^0=-\sqrt{\mathbf p^2+m^2}$ for particles in the final
state, and including the appropriate external line factors
$(2\pi)^{-3/2}u_\ell$ or $(2\pi)^{-3/2}v_\ell^*$ for initial particles or
antiparticles and $(2\pi)^{-3/2}u_\ell^*$ or $(2\pi)^{-3/2}v_\ell$ for final
particles or antiparticles. Indeed, when we come to the path integral approach
in Chapter 8 we shall find it easiest first to derive the Feynman\space%
